\documentclass[10pt]{article}
\usepackage{../../../estilo/cartoes}

\hypersetup{
  pdftitle={Vanguarda 44: Normandia - Cartoes de Cenario},
  pdfauthor={Felipe Ramires Terrazas},
}

\begin{document}

% ------------------------------------------------------------------
% Folha 1 -- Cartoes 1 e 2
% ------------------------------------------------------------------
\begin{cardpage}

\cardcell{0}{297}{1}{DIA-D: PRAIA DE OMAHA}{
\textbf{6 jun.\ 1944, Praia de Omaha, Normandia, França.}

\textbf{Forças:} Americanos (Atacante) vs.\ Alemães (Defensor).

\textbf{Qualidade de Tropa:} Alemães Regulares (guarnição costeira); Americanos Regulares.

\textbf{Regra Especial ``A Surpresa da Invasão'':} nos primeiros 2 turnos, unidades americanas que executarem Avançar ou Correr recebem +5cm extra de movimento, representando o avanço através da praia antes que a defesa alemã se reorganize por completo.

\textbf{Objetivo:} Americanos vencem se ocuparem a linha de bunkers até o fim do turno 8. Alemães vencem se os impedirem.

\textbf{Duração:} 8 turnos.
}

\cardcell{105}{297}{2}{CARENTAN}{
\textbf{10--13 jun.\ 1944, Carentan, Normandia, França.}

\textbf{Forças:} Americanos (Atacante) vs.\ Alemães (Defensor).

\textbf{Qualidade de Tropa:} ambos os lados Veteranos (paraquedistas de elite dos dois exércitos).

\textbf{Regra Especial ``Em Guarda'':} unidades alemãs que dispararem contra um alvo americano que executou Avançar ou Correr neste turno ganham +1 no Teste de Acerto: os defensores já esperavam o ataque pela estrada, com armas pré-posicionadas e comunicação estabelecida.

\textbf{Objetivo:} Americanos vencem se ocuparem o centro da cidade até o turno 8. Alemães vencem se resistirem.

\textbf{Duração:} 8 turnos.
}

\end{cardpage}

% ------------------------------------------------------------------
% Folha 2 -- Cartao 3
% ------------------------------------------------------------------
\begin{cardpage}

\cardcell{0}{297}{3}{CAEN}{
\textbf{Jul.\ 1944, Caen, Normandia, França.}

\textbf{Forças:} Britânicos (Atacante) vs.\ Alemães (Defensor).

\textbf{Qualidade de Tropa:} Alemães Veteranos (unidades blindadas endurecidas em combate); Britânicos Regulares.

\textbf{Regra Especial ``Bombardeio Prévio'':} no início da partida, role 1D6 para cada peça de terreno de edifício na mesa; num resultado de 5+, aquele edifício já está em ruínas: mantém Cobertura Pesada, mas também vira Terreno Difícil para atravessar.

\textbf{Objetivo:} Britânicos vencem se controlarem a maioria de 4 a 6 zonas de objetivo definidas na mesa até o turno 10.

\textbf{Duração:} 10 turnos.
}

\end{cardpage}

\end{document}
